\section{Reversible coupling and nonstationary transport}
\label{nuclear:10}
\subsection{Reversible extension of the balance between two forms}

Let \(B,C:V\to V'\) be isomorphisms preserving a nondegenerate form \(b\) between two realizations:

\[
B^*b'B=C^*b'C=b.
\]

The form may be a positive metric or an alternating form. On real spaces, \(*\) denotes transposition with respect to the chosen coordinates; the equality expresses transport of the form between the specified realizations. On Hilbert spaces, the operators and their inverses are bounded.

With \(s=\sqrt8\) and \(Q_V=\frac13\left(\begin{smallmatrix}sI&-I\\I&sI\end{smallmatrix}\right)\), define

\begin{equation}
\boxed{
\mathcal U(B,C)=Q_{V'}^*\operatorname{diag}(B,C)Q_V
=\frac19\begin{pmatrix}8B+C&\sqrt8(C-B)\\\sqrt8(C-B)&B+8C\end{pmatrix}.}
\label{nuc10:eq:1}
\end{equation}

Write \(T=(8B+C)/9\), \(D=\sqrt8(C-B)/9\), and \(R=(B+8C)/9\). The complete update is

\begin{equation}
\binom{x'}{m'}=\begin{pmatrix}T&D\\D&R\end{pmatrix}\binom{x}{m}.
\label{nuc10:eq:2}
\end{equation}

\begin{theorem}
The map \eqref{nuc10:eq:1} is reversible and preserves the form \(b\oplus b\):

\begin{equation}
\mathcal U(B,C)^*(b'\oplus b')\mathcal U(B,C)=b\oplus b,
\quad\mathcal U(B,C)^{-1}=\mathcal U(B^{-1},C^{-1}).
\label{nuc10:eq:3}
\end{equation}
\end{theorem}

\begin{proof}
We have \(Q^*Q=I\) because \(8+1=9\). Its scalar blocks make \(Q\) preserve \(b\oplus b\) for any form \(b\). The diagonal map preserves the forms by hypothesis. Composition proves \eqref{nuc10:eq:3}, and inversion of the three factors gives the declared inverse.
\end{proof}

For incoming memory \(m=0\), the balance \(b=T^*b'T+D^*b'D\) is recovered. The second column determines the response to arbitrary incoming memory. Conservation of positive norms and conservation of oriented area are the respective specializations of \eqref{nuc10:eq:3}.

\subsubsection{Collective realization in the nine copies}

On \(V^9\), let \(W_E=\operatorname{diag}(I,\ldots,I,E)\), with \(E\) in the ninth position. Let \(A(a,b)=(a/\sqrt8,\ldots,a/\sqrt8,b)\) and \(L=AQ\). Then \(L(x,0)=Jx\) and

\begin{equation}
W_E L=L\mathcal U(I,E).
\label{nuc10:eq:4}
\end{equation}

\begin{proof}
\(W_EA=A\operatorname{diag}(I,E)\); multiplying by \(Q\) and using \(AQ=L\) gives \eqref{nuc10:eq:4}. Orthogonality of the mean of the first eight positions to their seven differences also shows that the complement remains fixed under \(W_E\).
\end{proof}

The original cyclic step is \(S_E=P_{\rm cyc}W_E\). Hence \(S_EL=(P_{\rm cyc}L)U(I,E)\), with distinct input and output charts. The collective reduction preserves the corresponding pattern: \(S_E^9=I_9\otimes E\), whereas \(W_E^9=\operatorname{diag}(I_8,E^9)\). When a reader uses the nine individual labels, they are retained in their original representation.

\subsection{Exact composition and memory re-entry}

For \(B_1,C_1:V_0\to V_1\) and \(B_2,C_2:V_1\to V_2\),

\begin{equation}
\boxed{\mathcal U(B_2,C_2)\mathcal U(B_1,C_1)
=\mathcal U(B_2B_1,C_2C_1).}
\label{nuc10:eq:5}
\end{equation}

\begin{proof}
In \eqref{nuc10:eq:1}, the factors \(Q_{V_1}Q_{V_1}^*\) cancel, and the two diagonal blocks multiply in the indicated order. Induction extends the formula to any finite number of transports, preserving their order.
\end{proof}

The reading block contains the identity

\begin{equation}
\boxed{T_2T_1+D_2D_1=\frac{8B_2B_1+C_2C_1}{9}.}
\label{nuc10:eq:6}
\end{equation}

Starting from \((x,0)\), the first stage produces \((T_1x,D_1x)\). The second reading is \(T_2T_1x+D_2D_1x\). The second term is exactly the contribution of the reintroduced memory. Storing \(D_1x\) outside the next coupling and continuing only with \(T_1x\) gives the separate-register procedure, whose reading is \(T_2T_1x\). Equality \eqref{nuc10:eq:6} makes the difference between the two procedures precise.

The realization \(U\) preserves the coupling state and the ordered products \(B_N\cdots B_1\), \(C_N\cdots C_1\). Conservation of all prefixes additionally includes the HMT emission register, which retains each factor together with its order.

\subsubsection{Covariant increment equation}

For the separate register \(x_{n+1}=T_nx_n\), \(m_n=D_nx_n\), we have

\begin{equation}
x_{n+1}-B_nx_n=m_n/\sqrt8.
\label{nuc10:eq:7}
\end{equation}

With \(G_0=I\) and \(G_{n+1}=B_nG_n\), telescoping gives

\begin{equation}
x_0=G_N^{-1}x_N-\frac1{\sqrt8}\sum_{n<N}G_{n+1}^{-1}m_n.
\label{nuc10:eq:8}
\end{equation}

Memory is the increment relative to the transport between forms \(B_n\), with the declared normalization. Equality \eqref{nuc10:eq:8} is a finite algebraic reconstruction; its passage to an infinite series requires convergence. The following theorem constructs the stable inverse through the limit of positive operators.

\subsection{Memory of a nonstationary sequence}

In a Hilbert realization \(H\), let \(C_n\) be unitary, with their composition order declared. Write

\[
T_n=(8I+C_n)/9,\quad D_n=\sqrt8(C_n-I)/9,
\quad P_0=I,\quad P_{n+1}=T_nP_n,\quad A_N=P_N^*P_N.
\]

The index \(N\) counts compressions with separate storage, unlike complete transport with re-entry. For variable metrics with \(B_n\) isometric isomorphisms, conjugation by \(G_n\) takes the transports to a fixed Hilbert space and produces this same case, with \(\widetilde C_n=G_{n+1}^{-1}C_nG_n\). The assertion uses the fibre metrics and the declared isometries between them.

\begin{theorem}[Positive terminal component and nonstationary recovery]
\label{nuclear:terminal-no-estacionario}
There exists a positive operator \(0\leq A_\infty\leq I\), the strong limit of \(A_N\), and

\begin{equation}
\boxed{I=A_\infty+\sum_{n\ge0}P_n^*D_n^*D_nP_n.}
\label{nuc10:eq:9}
\end{equation}

The map

\begin{equation}
\boxed{\mathcal Vv=(A_\infty^{1/2}v,(D_nP_nv)_{n\ge0})}
\label{nuc10:eq:10}
\end{equation}

is isometric and has a stable left inverse

\begin{equation}
\mathcal V^*(a,(m_n))=A_\infty^{1/2}a+
\sum_{n\ge0}P_n^*D_n^*m_n,\qquad\|\mathcal V^*\|=1.
\label{nuc10:eq:11}
\end{equation}

\end{theorem}
\begin{proof}
The local identity \(T_n^*T_n+D_n^*D_n=I\) gives \(A_N-A_{N+1}=P_N^*D_N^*D_NP_N\geq0\). The quadratic forms of the decreasing positive sequence converge to that of a positive operator \(A_\infty\). For \(B_N=A_N-A_\infty\), the inequality \(0\leq B_N^2\leq B_N\) makes this convergence strong. Finite telescoping and the limit prove \eqref{nuc10:eq:9}. Evaluation on \(v\) proves \eqref{nuc10:eq:10}. Its adjoint is \eqref{nuc10:eq:11}; the series converges in norm because truncations of a square-summable register converge and the adjoint is bounded.
\end{proof}

Reconstruction with the limiting terminal component and \(N\) registers has exact error \((A_N-A_\infty)v\). The vector \(A_\infty^{1/2}v\) is the canonical positive realization of the norm terminal component. Its existence follows from the Gram limit, independently of vector convergence of \(P_Nv\).

\subsection{Three distinct properties: exhaustion, distinguishability, and stability}

Let \(Mv=(D_nP_nv)_n\). By \eqref{nuc10:eq:9}, \(M^*M=I-A_\infty\). Therefore:

\begin{enumerate}
\item The entire norm of \(v\) passes to the register exactly when \(A_\infty v=0\).
\item The register distinguishes every state exactly when \(\ker(I-A_\infty)=\{0\}\).
\item Recovery from memory alone has a bounded inverse exactly when \(I-A_\infty\geq\varepsilon I\) for some \(\varepsilon>0\).
In that case, a left inverse is \((I-A_\infty)^{-1}M^*\).
\end{enumerate}

Moreover,

\begin{equation}
\boxed{\ker M=\bigcap_{n\ge0}\operatorname{Fix}(C_n).}
\label{nuc10:eq:12}
\end{equation}

\begin{proof}
If \(Mv=0\), then \(D_0v=0\), whence \(C_0v=v\) and \(P_1v=v\). Inductively, \(D_nP_nv=0\) gives \(C_nv=v\) and \(P_{n+1}v=v\). The converse is immediate. The three equivalences follow from the Gram matrix \(M^*M\) and the characterization of operators bounded below.
\end{proof}

\subsubsection{An exact example separating the three properties}

In the algebraic example \(C_0=-I\) and \(C_n=I\) for \(n\geq1\),

\begin{equation}
A_\infty=\frac{49}{81}I,\qquad M^*M=\frac{32}{81}I,
\qquad m_0=-\frac{2\sqrt8}{9}v,
\qquad v=-\frac9{2\sqrt8}m_0.
\label{nuc10:eq:13}
\end{equation}

Memory contains \(32/81\) of the squared norm and makes it possible to recover the complete vector. The terminal component retains \(49/81\) and is a positive operator distinct from a projector. Recoverability depends on the kernel and conditioning of the reader. Reading and memory constitute correlated components of an isometry.

This example separates the three logical consequences of the theorem. Application to a physical trajectory additionally preserves the selection and effective origin of its \(C_n\) through the TPK.

\subsection{Quantitative criterion for complete transfer}

For \(P_{j,n}=T_{j-1}\cdots T_n\) and \(P_{n,n}=I\), a window of length \(L\) has Gram matrix

\[
G_{n,L}=\sum_{j=n}^{n+L-1}P_{j,n}^*D_j^*D_jP_{j,n}
=I-P_{n+L,n}^*P_{n+L,n}.
\]

If \(G_{n,L}\geq\kappa I\), with \(\kappa>0\), in each consecutive block of length \(L\), then

\begin{equation}
\|P_N\|^2\le(1-\kappa)^{\lfloor N/L\rfloor},\qquad A_\infty=0.
\label{nuc10:eq:14}
\end{equation}

Each block contracts the squared norm by \(1-\kappa\); composition and contractivity of the remaining steps prove \eqref{nuc10:eq:14}.

A sufficient criterion preceding the products is

\begin{equation}
\sum_{j=n}^{n+L-1}D_j^*D_j\ge\alpha I
\quad\Longrightarrow\quad
G_{n,L}\ge\frac{\alpha}{[1+2(L-1)/9]^2}I.
\label{nuc10:eq:15}
\end{equation}

\begin{proof}
For \(y_j=P_{j,n}v\), we have \(y_{j+1}-y_j=D_jy_j/\sqrt8\) and \(\|D_j\|\leq2\sqrt8/9\). Hence
\[
\|D_jv\|\leq\|D_jy_j\|+\frac29\sum_{i=n}^{j-1}\|D_iy_i\|.
\]
The triangular matrix of partial sums has norm at most \(L-1\). Applying the window's \(\ell^2\) norm and the Gram hypothesis proves \eqref{nuc10:eq:15}.

\end{proof}

The absence of a common fixed vector gives distinguishability, whereas \eqref{nuc10:eq:15} adds a quantitative uniform separation. Neither condition is presumed merely by mentioning symmetries: each applies to the effective transports of the sequence.

\subsection{Complete incidence realization and changes of regime}

Section~\ref{nuc05:isometria} constructs, for the twelve-coordinate register, \(Fk=(\mu(k),\mathcal IP_0k/6)\), with an inverse and weighted inner product on \(Y=\RR\oplus\mathcal I(\mathbf1^\perp)\). The transformation \(F\oplus F\) takes the complete coupling to

\begin{equation}
\mathcal U_Y=(F'\oplus F')\mathcal U(B,C)(F^{-1}\oplus F^{-1})
=\mathcal U(F'BF^{-1},F'CF^{-1}).
\label{nuc10:eq:16}
\end{equation}

This identity preserves the second column, re-entry, and composition. For variable forms, the form is also transported through \(F\). In the real realization of the twelve \(A_2\) fibres, use \(F\otimes\operatorname{id}_{V_{A_2}}\), with \(V_{A_2}=A_2\otimes_{\ZZ}\RR\). Its domain is the real space of fibres; preservation of the integer lattice corresponds to the lattice transport already proved. The map \(F\) recovers \(K\) and retains the remaining genealogical labels in the complete state.

The positive version of \eqref{nuc10:eq:9}–\eqref{nuc10:eq:15} uses unitary transports in the declared metrics. The alternating-form version of \eqref{nuc10:eq:1}–\eqref{nuc10:eq:8} admits the three Euler regimes. With \(E=F_{\rm av}^2=\left(\begin{smallmatrix}1&1\\1&2\end{smallmatrix}\right)\), we obtain \(\det T=89/81\), \(\det D=-8/81\). The finite oriented sums telescope, and the area terminal components grow as \((89/81)^N\). The positive Hilbert limit uses the specific unitary hypotheses of Theorem~\ref{nuclear:terminal-no-estacionario}.

The balance brings together forms and domains, reversible evolution with re-entry, and unlimited recovery in the positive realizations. The subsequent physical applications preserve the readers, connections, and scales of this operator composition.
